\scanpage{62}
\mspage{58}
% archive sheet number: 61

On trouve une catégorie, \deleted{la catégorie} $\unsure{\mathrm{ER}|E\mathcal R|E\mathscr R}$, et ses objets (\unsure{équival\textsuperscript{t}} : des graphes (stricts)), \unsure{mériteraient} un nom, exprimant l'intuition qui s'attache à la situation.
On pourrait appeler un $(E,\unsure{|\mathrm{o}|_E|\lvert\circ\rvert_E})$ un ensemble avec relation de disjonction, un \uline{disjoncteur}, et \deleted{la catégorie \unsure{précéd\textsuperscript{te}}} parler des \uline{morphismes de disjoncteurs}, notation $\unsure{\operatorname{Hom}^{\mathrm{dis}}|\operatorname{Homdis}|\operatorname{Hom}_{\mathrm{dis}}}(\underline E,\underline E')$, et la \uline{catégorie des disjoncteurs}.
Ceci dit, \deleted{\unsure{l'on}} pour un \inserted{tel} morphisme de disjoncteurs
\[
E\xrightarrow{\alpha}\unsure{\mathcal P|\mathscr P|P}(E')
\]
considérons l'application correspondante
\[
\deleted{\unsure{\mathbb\Sigma|\Sigma}_{\unsure{E}}}
\]
\diagram{
\begin{tikzcd}[ampersand replacement=\&,column sep=large]
\overline\alpha:\ \mathcal P(E) \arrow[r] \& \mathcal P(E') \\
A \arrow[r,mapsto] \& \displaystyle\bigcup_{x\in A}\alpha(x)
\end{tikzcd}}
et l'application composée
\diagram{
\begin{tikzcd}[ampersand replacement=\&,column sep=large]
\mathbb\Sigma_{\unsure{\underline E|E}} \arrow[r,hook,"\mathrm{inc}"] \& \mathcal P(E) \arrow[r,"\overline\alpha"] \& \mathcal P(E') \arrow[r,"\operatorname{supp}^{\unsure{\mathrm{o}|\circ|0}}_{E'}"] \& \mathbb\Sigma_{E'}
\end{tikzcd}}
,
\deleted{Il est} on pose
\[
\alpha_*:\mathbb\Sigma_E\longrightarrow\mathbb\Sigma_{E'},\qquad
\alpha_{\unsure{\circ|*}}(S)=\operatorname{supp}^{\mathrm{o}}_{E'}(\overline\alpha(S)).
\]
Cette application est en fait définie \deleted{\unsure{pour}}\inserted{\unsure{quelque|quelle}}
